\chapter{Fleets and Ships}

\section{Aurelian Compact}

Heirs to the old gate authority, the Compact offers reunification at lance-point. Aurelian ships reward deliberate positioning, overlapping beam arcs, and composed formation fighting.

\begin{tabularx}{\linewidth}{l c c c c Y}\toprule
Ship & Size & Energy & Hull & Shields F/A & Weapons\\\midrule
Aurelian Frigate & Small & 8 & 5 & 4/4 & 4 beams, 1 driver\\
Aurelian Cruiser & Medium & 11 & 7 & 6/5 & 4 beams, 2 drivers\\
Aurelian Battleship & Large & 15 & 9 & 9/7 & 8 beams, 2 drivers\\\bottomrule
\end{tabularx}

\begin{description}[style=nextline,leftmargin=0pt]
  \item[Aurelian Frigate] Swift Compact escorts race for the lane mouth, then rake intruders with disciplined crossing fire.
  \item[Aurelian Cruiser] Compact cruisers broadcast terms of reunification while paired beam batteries make refusal expensive.
  \item[Aurelian Battleship] A mobile fragment of the lost capital, built to make distant worlds remember Aurelian authority.
\end{description}

\section{Veyr Dominion}

A fortress-state forged by isolation, the Dominion trusts centralized command and decisive force. Veyr ships combine specialized batteries with missiles and long-range pressure.

\begin{tabularx}{\linewidth}{l c c c c Y}\toprule
Ship & Size & Energy & Hull & Shields F/A & Weapons\\\midrule
Veyr Frigate & Small & 7 & 5 & 3/3 & 2 beams, 1 driver, 1 launcher\\
Veyr Cruiser & Medium & 10 & 7 & 5/5 & 4 beams, 1 driver, 1 launcher\\
Veyr Battleship & Large & 14 & 8 & 8/7 & 2 drivers, 1 launcher\\\bottomrule
\end{tabularx}

\begin{description}[style=nextline,leftmargin=0pt]
  \item[Veyr Frigate] A patient hunter that marks targets for the Dominion and closes only when escape is impossible.
  \item[Veyr Cruiser] The workhorse of Veyr blockade lines, equally prepared to hold a corridor or break one open.
  \item[Veyr Battleship] The Dominion's answer to uncertainty: armored command, long-range fire, and no avenue of retreat.
\end{description}

\section{Kestrel Freeholds}

Station clans and convoy captains fight to keep the reopened frontier unowned. Their vessels favor resilient shields, flexible missile coverage, and hulls adapted by independent yards.

\begin{tabularx}{\linewidth}{l c c c c Y}\toprule
Ship & Size & Energy & Hull & Shields F/A & Weapons\\\midrule
Kestrel Frigate & Small & 9 & 4 & 4/5 & 2 beams, 2 launchers\\
Kestrel Cruiser & Medium & 12 & 6 & 6/6 & 2 beams, 3 launchers\\
Kestrel Battleship & Large & 16 & 8 & 9/8 & 6 beams, 3 launchers\\\bottomrule
\end{tabularx}

\begin{description}[style=nextline,leftmargin=0pt]
  \item[Kestrel Frigate] Courier frames turned raiders, fast enough to strike a claim beacon and vanish into the fracture lanes.
  \item[Kestrel Cruiser] Freehold yards wrap stubborn shields around missile magazines built to outlast wealthier enemies.
  \item[Kestrel Battleship] No two are identical; each is a shipyard covenant made armor, carrying a whole Freehold to war.
\end{description}

\chapter{First Light Tutorial}

The browser prototype provides a guided First Light scenario. For a tabletop version, use one Aurelian frigate and one Kestrel frigate on a 12 by 12 map, 10 hexes apart and facing each other.

\begin{enumerate}
  \item \term{Allocation}: choose a modest speed and charge one lance beam. Leave enough energy to reinforce the shield facing the enemy.
  \item \term{Movement}: announce speed, roll initiative, and draw the first impulse card. Move only a ship whose speed is printed on the card.
  \item \term{Turn mode and side-slip}: track translated hexes, then demonstrate one legal turn and one side-slip. Remember that a second consecutive side-slip is forbidden.
  \item \term{Missiles}: launch after movement. Leave the missile in the launch hex until the next impulse, then move it two hexes toward its target.
  \item \term{Direct fire}: choose a charged weapon, verify arc and range, then roll against the weapon matrix.
  \item \term{Shields and damage}: apply reinforcement, then shield boxes, then roll once for each penetrating damage point.
  \item \term{End the turn}: repair a shield if funded, clear allocations, and begin again.
\end{enumerate}
